\chapter[CAPACITÉS, REQUÊTES ET RÉPONSES DE STATION]{Capacités, requêtes et réponses de
station}\label{capacituxe9s-requuxeates-et-ruxe9ponses-de-station}

\section{Capacités de station}\label{capacituxe9s-de-station}

Une station peut définir un ensemble d'attributs --- les Station
Capabilities --- transmis en réponse à une requête \texttt{IGATE}, avec
DTI \texttt{\textless{}}. Chaque capacité est un \texttt{TOKEN} ou une
paire \texttt{TOKEN=VALUE}. Plusieurs sur une ligne, séparées par des
virgules.

Capacités actuellement définies :

\begin{verbatim}
IGATE,MSG_CNT=n,LOC_CNT=n
\end{verbatim}

\begin{itemize}
\tightlist
\item
  \texttt{IGATE} : la station est un IGate.
\item
  \texttt{MSG\_CNT} : nombre de messages transmis.
\item
  \texttt{LOC\_CNT} : nombre de stations « locales » (à qui l'IGate
  transmet les messages sur le RF local).
\end{itemize}

\section{Requêtes et réponses}\label{requuxeates-et-ruxe9ponses}

Deux types : générale (à toutes les stations) ou dirigée (à une station
individuelle en format message).

Les requêtes commencent toujours par \texttt{?}, sont envoyées une seule
fois, sans identifiant de message et ne sont pas acquittées. Idem pour
les réponses.

Types de requête :

\begin{longtable}[]{@{}
  >{\raggedright\arraybackslash}p{(\columnwidth - 4\tabcolsep) * \real{0.3333}}
  >{\raggedright\arraybackslash}p{(\columnwidth - 4\tabcolsep) * \real{0.3333}}
  >{\raggedright\arraybackslash}p{(\columnwidth - 4\tabcolsep) * \real{0.3333}}@{}}
\toprule\noalign{}
\begin{minipage}[b]{\linewidth}\raggedright
Type
\end{minipage} & \begin{minipage}[b]{\linewidth}\raggedright
Portée
\end{minipage} & \begin{minipage}[b]{\linewidth}\raggedright
Réponse
\end{minipage} \\
\midrule\noalign{}
\endhead
\bottomrule\noalign{}
\endlastfoot
\texttt{APRS} & Générale & Position + statut de la station \\
\texttt{APRSD} & Dirigée --- stations entendues directement & Liste des
stations entendues direct \\
\texttt{APRSH} & Dirigée --- as-tu entendu telle station ? & Position de
la station entendue comme Object APRS + stats des 8 dernières heures \\
\texttt{APRSM} & Dirigée --- messages non ACK/livrés en attente & Tous
les messages en suspens \\
\texttt{APRSO} & Dirigée --- Objects & Objects de la station \\
\texttt{APRSP} & Dirigée --- position & Position \\
\texttt{APRSS} & Dirigée --- statut & Statut \\
\texttt{APRST} ou \texttt{PING?} & Dirigée --- trace & Route trace \\
\texttt{IGATE} & Générale --- tous les IGates & Capacités IGate \\
\lit{WX} & Générale --- toutes les stations météo & Weather report (+
position si absente) \\
\end{longtable}

Une station interrogée sans info à retourner peut ignorer la requête.
Toute requête non reconnue est ignorée.

\section{Requêtes générales}\label{requuxeates-guxe9nuxe9rales}

\begin{center}
\begin{tabular}{|c|c|c|c|}
\hline
\textbf{Champ} & \lit{?} & Query Type & \lit{?} \\
\hline
\textbf{Octets} & 1 & n & 1 \\
\hline
\end{tabular}
\end{center}
\noindent + Target Footprint optionnel : \texttt{Lat, Long, Radius} (lat/long en degrés flottants, radius en miles sur 4 chiffres).

Exemples :

\begin{longtable}[]{@{}
  >{\raggedright\arraybackslash}p{(\columnwidth - 2\tabcolsep) * \real{0.5000}}
  >{\raggedright\arraybackslash}p{(\columnwidth - 2\tabcolsep) * \real{0.5000}}@{}}
\toprule\noalign{}
\begin{minipage}[b]{\linewidth}\raggedright
Requête
\end{minipage} & \begin{minipage}[b]{\linewidth}\raggedright
Réponse typique
\end{minipage} \\
\midrule\noalign{}
\endhead
\bottomrule\noalign{}
\endlastfoot
\texttt{?APRS?} & Position + statut standard \\
\texttt{?APRS?\ 34.02,-117.15,0200} & Cible : rayon 200 miles autour de
34.02°N, 117.15°W \\
\texttt{?IGATE?} & \texttt{\textless{}IGATE,MSG\_CNT=43,LOC\_CNT=14} \\
\texttt{?WX?} & Weather Report standard + position \\
\end{longtable}

Pour \texttt{?APRS?} avec footprint, lat/long en degrés flottants (pas
en format APRS). N/E positifs → espace initial. S/W négatifs. Rayon en
miles sur 4 chiffres fixes.

\section{Requêtes dirigées}\label{requuxeates-diriguxe9es}

Format message APRS sans identifiant. Addressee = indicatif de la
station cible.

\begin{center}
\begin{tabular}{|c|c|c|c|c|c|}
\hline
\textbf{Champ} & \lit{:} & Addressee & \lit{:?} & Query Type & Callsign of Heard Station \\
\hline
\textbf{Octets} & 1 & 9 & 2 & 5 & 0-9 (optionnel) \\
\hline
\end{tabular}
\end{center}

Exemples : - \texttt{:KH2Z\ \ \ \ \ :?APRSD} →
\texttt{:N8UR\ \ \ \ \ :Directs=\ WA1LOU\ WD5IV\ D…} -
\texttt{:KH2Z\ \ \ \ \ :?APRSH\ N0QBF} →
\texttt{:N8UR\ \ \ \ \ :N0QBF\ HEARD:\ 1\ 3\ 2\ .\ .\ 4\ 5\ 6} -
\texttt{:KH2Z\ \ \ \ \ :?APRSM} → tous messages non ACK/livrés. -
\texttt{:KH2Z\ \ \ \ \ :?APRSO} → objects. -
\texttt{:KH2Z\ \ \ \ \ :?APRSP} → position. -
\texttt{:KH2Z\ \ \ \ \ :?APRSS} → statut. -
\texttt{:KH2Z\ \ \ \ \ :?APRST} ou \texttt{?PING?} → trace.

\begin{center}\rule{0.5\linewidth}{0.5pt}\end{center}
